
\section{Algorithm}

\begin{algorithm}
\caption{gPCD Broad-Phase and Narrow-Phase Collision Detection}
\begin{algorithmic}[1]

\Require Particle array $P$ of length $N$
\Require Uniform cell array $C$
\Require Grid dimensions $N_x, N_y, N_z$
\Ensure Collision set $K$

\State Construct cell array $C$ with cell size $1 \times 1 \times 1$
\State Initialize all occupancy lists in $C$ as empty
\State Initialize collision set $K \gets \emptyset$

\Statex
\Statex \textbf{Broad Phase: Graphics Pipeline}

\ForAll{particles $p$ in parallel at vertex stage}
    \State Emit $8$ corner fragment invocations
\EndFor

\ForAll{corner fragments of particle $p$ in parallel}
    \State $c \gets$ unique corner ID in $\{0,1,\dots,7\}$
    \State $\mathbf{x}_c \gets$ corner position of particle $p$
    \State $(i,j,k) \gets \lfloor \mathbf{x}_c \rfloor$
    \State $I \gets \Call{Flatten}{i,j,k}$
    \State $P[p].\textit{cornerLocation}[c] \gets I$
    \State Insert particle $p$ into occupancy list of cell $C[I]$
\EndFor

\Statex
\Statex \textbf{Narrow Phase: Compute Pipeline}

\ForAll{particles $p$ in parallel at compute stage}
    \State Local duplicate array $D \gets \emptyset$

    \For{$c = 0$ to $7$}
        \State $I \gets P[p].\textit{cornerLocation}[c]$

        \ForAll{candidate particles $q$ in occupancy list of cell $C[I]$}
            \If{$q = p$}
                \State \textbf{continue}
            \EndIf

            \If{$q \in D$}
                \State \textbf{continue}
            \EndIf

            \State Insert $q$ into $D$

            \If{$\|\mathbf{x}_p - \mathbf{x}_q\| < r_p + r_q$}
                \State Insert pair $(p,q)$ into $K$
            \EndIf
        \EndFor
    \EndFor
\EndFor

\State \Return $K$

\end{algorithmic}
\end{algorithm}

\begin{algorithm}
\caption{Flattening a 3D Cell Coordinate to a 1D Index}
\begin{algorithmic}[1]

\Function{Flatten}{$i,j,k$}
    \State \Return $i + N_x \left( j + N_y k \right)$
\EndFunction

\end{algorithmic}
\end{algorithm}

\subsection{Corner Assignment Shader (GLSL-style)}

The following pseudocode illustrates the corner-based cell assignment as implemented in a shader-compatible form. Each invocation corresponds to a single particle corner and computes its associated cell index.

\begin{lstlisting}[language=C++, caption={Corner-to-cell mapping in GLSL/Vulkan-style pseudocode}]
shared int cornerCell[8];

void assignCornerCell(uint corner, vec3 pos, float radius)
{
    // Decode corner index into sign vector using bitwise operations
    vec3 sign = vec3(
        ((corner & 1u) != 0u) ?  1.0 : -1.0,
        ((corner & 2u) != 0u) ?  1.0 : -1.0,
        ((corner & 4u) != 0u) ?  1.0 : -1.0
    );

    // Compute corner position of AARB
    vec3 cornerPos = pos + sign * radius;

    // Determine containing cell
    ivec3 cell = ivec3(floor(cornerPos));

    // Convert 3D cell coordinate to linear index
    int cellID = flatten(cell);

    // Store result in per-particle corner array
    cornerCell[corner] = cellID;
}
\end{lstlisting}

\paragraph{Implementation Notes}

The corner index is interpreted as a 3-bit mask, where each bit determines the sign of the offset in one spatial dimension. This formulation generates all eight corners of the particle’s axis-aligned bounding region without branching or lookup tables.

Each shader invocation writes to a unique index in the \texttt{cornerCell} array, ensuring write-conflict-free execution under parallel scheduling. The use of direct index computation enables constant-time cell assignment and predictable memory access patterns.